\documentclass[../DoAn.tex]{subfiles}
\begin{document}

Chương trước đã đặc tả đầy đủ các yêu cầu chức năng và phi chức năng của hệ
thống. Để biến các đặc tả đó thành một sản phẩm vận hành được, đề tài lựa
chọn một bộ công nghệ với tiêu chí tổng thể là \textit{hiệu năng cao trên đường
nhanh, có khả năng tích hợp với hệ sinh thái học máy, và triển khai được trong
môi trường mã nguồn mở khép kín}. Chương này trình bày nền tảng lý thuyết và
đặc tính kỹ thuật của các thành phần công nghệ được lựa chọn, kèm theo lý giải
vì sao mỗi lựa chọn phù hợp với yêu cầu đã xác định ở Chương \ref{section:2.4}.
Các thành phần được tổ chức theo lớp kiến trúc: từ ngôn ngữ và runtime của
WAF core (Go), đến cơ chế phát hiện (Rule Engine, DistilBERT, FastAPI), hạ
tầng lưu trữ (PostgreSQL), lớp bảo mật nội bộ (JWT, Bcrypt, RBAC), thuật toán
điều khiển tải (Token Bucket), và cuối cùng là môi trường đóng gói triển khai
(Docker, Docker Compose).

% ====================================================================
\section{Go (Golang) -- Nền tảng xây dựng WAF core}
\label{section:3.1}

\subsection{Khái quát}
\label{subsection:3.1.1}

Go là ngôn ngữ lập trình biên dịch tĩnh được Google công bố năm 2009, hướng
đến giải quyết các vấn đề thực tế phát sinh khi xây dựng phần mềm hệ thống
quy mô lớn. Ba đặc tính cốt lõi của Go bao gồm: \textbf{(i)} mô hình đồng thời
(concurrency model) dựa trên \textit{goroutine} -- các luồng xanh nhẹ hơn
luồng hệ điều hành nhiều lần (mặc định stack 8 KB so với hàng MB của thread
truyền thống), được lập lịch bởi runtime của ngôn ngữ theo mô hình M:N (M
goroutine ánh xạ lên N thread của OS); \textbf{(ii)} cơ chế giao tiếp giữa các
goroutine qua \textit{channel} thay vì shared memory với mutex như mô hình
cổ điển, được tổng kết bởi phương châm ``don't communicate by sharing memory;
share memory by communicating''; và \textbf{(iii)} thư viện chuẩn (standard
library) phong phú và ổn định, bao gồm gói \texttt{net/http} có thể xây dựng
máy chủ HTTP đầy đủ tính năng mà không cần bất kỳ framework bên ngoài nào.

Về quản lý bộ nhớ, Go sử dụng garbage collector concurrent tri-color
mark-and-sweep từ phiên bản 1.5, với mục tiêu duy trì thời gian dừng
(\textit{stop-the-world pause}) dưới một mili giây ngay cả với heap kích
thước hàng GB. Đây là đặc tính quan trọng đối với các ứng dụng yêu cầu độ
trễ thấp và tính xác định cao như reverse proxy.

\subsection{Vai trò trong hệ thống}
\label{subsection:3.1.2}

Toàn bộ WAF core trong đề tài được viết bằng Go: vòng lặp chấp nhận kết nối
HTTP, parser request, normalizer, rule engine, các middleware (xác thực, rate
limit, logging), bộ điều phối quyết định, và lớp adapter giao tiếp với ML
service. Riêng phần ML inference được tách ra microservice Python (mục
\ref{section:3.4}) để tận dụng hệ sinh thái Hugging Face -- tuy nhiên việc
giao tiếp giữa Go và service Python được thực hiện qua HTTP với kiểu dữ liệu
đơn giản, không tạo ràng buộc chặt giữa hai runtime.

Gói \texttt{net/http} của thư viện chuẩn cung cấp interface \texttt{Handler}
và pipeline middleware dạng chain -- chính là khung mà WAF middleware stack
được xây dựng. Đặc biệt, \texttt{httputil.ReverseProxy} cung cấp sẵn cơ chế
tunnel ngược cho HTTP/1.1, HTTP/2 và WebSocket, giúp đề tài không phải hiện
thực hóa lại logic forwarding ở mức byte. Việc dispatching theo prefix path
được thực hiện qua \texttt{http.ServeMux} với mẫu kết thúc bằng dấu ``/'' để
khớp tiền tố -- mẫu này tương thích với phiên bản Go 1.21 mà đề tài cố định
nhằm giảm phụ thuộc vào các tính năng mới (như pattern matching
\texttt{\{id\}} chỉ có từ Go 1.22).

\subsection{Lý do lựa chọn}
\label{subsection:3.1.3}

Việc lựa chọn Go cho lớp WAF core được dựa trên ba so sánh thực tế với các
ứng cử viên thường gặp.

So với \textbf{Python}, Go nhanh hơn về throughput I/O-bound từ 5 đến 20 lần
trong các benchmark độc lập (TechEmpower Round 22, 2023), nhờ goroutine
không phải chịu chi phí của Global Interpreter Lock (GIL). Hơn nữa, biên dịch
ra binary tĩnh đơn lẻ giúp triển khai gọn nhẹ, không phụ thuộc môi trường
Python ở máy đích -- đặc tính quan trọng cho image Docker.

So với \textbf{Node.js}, Go cung cấp hiệu năng tương đương ở các workload
I/O-bound đơn giản, nhưng vượt trội rõ rệt khi workload tính toán hỗn hợp
(parser, regex matching, mã hóa) do mô hình đa luồng thật thay vì event loop
đơn luồng. Việc xử lý song song nhiều request mà không phải lập trình
asynchronous tường minh giúp mã nguồn dễ đọc và dễ kiểm thử hơn.

So với \textbf{Rust}, Go thua kém về hiệu năng tuyệt đối và an toàn bộ nhớ
ở mức compile-time, nhưng vượt trội về tốc độ phát triển và độ thân thiện
với người mới. Rust phù hợp hơn cho các thành phần ở tầng thấp như packet
parser ở Layer 3/4, vốn không phải phạm vi của đề tài. Đối với một reverse
proxy ở tầng 7 với độ trễ mục tiêu vài mili giây, Go cân bằng tốt nhất giữa
ba tiêu chí: hiệu năng đủ dùng, tốc độ phát triển nhanh, và hệ sinh thái
phong phú.

% ====================================================================
\section{Rule Engine và OWASP CRS -- Cơ sở phát hiện dựa trên luật}
\label{section:3.2}

\subsection{Khái quát cơ chế rule-based detection}
\label{subsection:3.2.1}

Rule-based detection là mô hình phát hiện tấn công cổ điển nhất trong không
gian WAF, dựa trên ý tưởng nền tảng: \textit{nếu một request có chứa các mẫu
văn bản đặc trưng của một loại tấn công đã biết, thì request đó nhiều khả
năng là tấn công}. Mô hình này được hiện thực hóa bằng tập rule -- mỗi rule
mô tả một mẫu pattern cần kiểm tra trên một hay nhiều trường của request,
kèm theo hành động cần thực hiện khi pattern khớp.

Hệ thống đề tài kế thừa và tinh giản mô hình \textit{anomaly scoring} mà
OWASP Core Rule Set đã đề xuất từ phiên bản 3.0. Khác với mô hình
\textit{immediate block} (chặn ngay khi có rule khớp đầu tiên), anomaly
scoring cho phép mỗi rule đóng góp một điểm số nhỏ vào tổng điểm nguy cơ
của request; quyết định block chỉ được đưa ra khi điểm tổng vượt một ngưỡng
được cấu hình. Ưu điểm chính của mô hình này là giảm false positive cho các
trường hợp một request bình thường vô tình khớp một rule yếu (ví dụ chứa từ
khóa SQL trong dữ liệu hợp lệ), nhưng vẫn giữ được khả năng phát hiện các
tấn công đa pha mà mỗi pattern đơn lẻ không đủ kết luận.

Một trong những thách thức cơ bản của rule-based detection là \textit{kỹ
thuật vượt mặt} (evasion). Kẻ tấn công có thể biến tấu payload thông qua
URL encoding nhiều lớp, thay đổi hoa thường, chèn ký tự trắng (whitespace
padding), hay sử dụng biểu thức tương đương ngữ nghĩa nhưng khác cú pháp.
Để chống lại các kỹ thuật này, mỗi rule được trang bị một \textit{transform
chain} -- chuỗi các biến đổi tiền xử lý áp dụng trên text đầu vào trước khi
so khớp pattern. Chuỗi mặc định bao gồm:

\begin{itemize}[leftmargin=*]
    \item \texttt{URL\_DECODE}: giải mã URL một lớp, chuyển \texttt{\%XX}
    về ký tự tương ứng -- chống obfuscation cơ bản.
    \item \texttt{LOWERCASE}: chuyển toàn bộ về chữ thường -- chống
    case-mixing kiểu \texttt{uNiOn SeLeCt}.
    \item \texttt{COMPRESS\_WHITESPACE}: dồn các khoảng trắng liên tiếp về
    một ký tự duy nhất -- chống chèn ký tự trắng giữa các từ khóa.
    \item \texttt{HTML\_ENTITY\_DECODE}: giải mã các entity HTML
    (\texttt{\&lt;}, \texttt{\&\#x3c;}) -- quan trọng với các XSS payload.
    \item \texttt{BASE64\_DECODE}: hỗ trợ giải mã base64 cho các nhóm rule
    đặc thù như Log4Shell với JNDI URL được encode base64.
\end{itemize}

Đối với pattern matching, đề tài hỗ trợ hai kiểu mẫu. \textbf{REGEX pattern}
dùng biểu thức chính quy của Go (\texttt{regexp} package, dựa trên RE2,
không hỗ trợ backreference nhưng đảm bảo độ phức tạp tuyến tính, không bị
catastrophic backtracking). \textbf{TOKEN pattern} dựa trên khái niệm
\textit{proximity match}: liệt kê một dãy token và yêu cầu chúng xuất hiện
trong văn bản theo đúng thứ tự, trong khoảng cách tối đa N ký tự. Token
pattern thường mạnh hơn regex trong việc phát hiện các tấn công đa từ khóa
mà thứ tự quan trọng nhưng khoảng đệm có thể thay đổi (ví dụ
\texttt{union}\ldots\texttt{select} cách nhau bởi comment SQL).

\subsection{Cấu trúc rule JSON}
\label{subsection:3.2.2}

Mỗi rule được mô tả dưới dạng một đối tượng JSON với các trường chính: định
danh và metadata (\texttt{id}, \texttt{category}, \texttt{severity},
\texttt{description}), điều kiện áp dụng (\texttt{conditions} -- các trường
của request mà rule cần kiểm tra: \texttt{path}, \texttt{query},
\texttt{headers.user-agent}, \texttt{body}, v.v.), chuỗi biến đổi tiền xử lý
(\texttt{transforms}), danh sách pattern (\texttt{patterns} -- có thể là
mảng các regex hoặc token sequence), thông số scoring (\texttt{anomaly\_score},
\texttt{severity\_multiplier}), hành động (\texttt{actions} -- mặc định là
\textit{score-only}, có thể nâng lên \textit{deny-immediate} cho các rule
critical), và danh sách ngoại lệ (\texttt{exceptions} -- ví dụ bỏ qua đường
dẫn \texttt{/health}, \texttt{/metrics}).

Cấu trúc JSON có hai ưu điểm chính so với cú pháp ngôn ngữ chuyên dụng của
ModSecurity (SecRule DSL). Thứ nhất, JSON là chuẩn dữ liệu phổ biến, có thể
được sinh tự động từ công cụ và validate bằng JSON Schema. Thứ hai, việc
nạp lại tập rule mà không cần khởi động lại proxy được thực hiện đơn giản
qua một endpoint admin nhận file JSON, parse, validate, và hot-swap engine
state -- hành vi mà ModSecurity vốn không hỗ trợ tốt do kiến trúc module
nhúng vào web server.

Mỗi phase của vòng đời request có thể có rule riêng. Trong phạm vi đề tài,
chỉ phase REQUEST được hiện thực hóa -- kiểm tra request trước khi forward
đến backend. Phase RESPONSE (kiểm tra response trước khi trả về cho client,
phát hiện data leak hay error message) được xác định là hướng phát triển
tương lai (mục \ref{section:6.2}).

\subsection{Bộ luật 78 rules / 13 nhóm}
\label{subsection:3.2.3}

Đề tài xây dựng tập luật gồm 78 rule chia thành 13 nhóm tấn công, bao phủ
hầu hết các nhóm trong OWASP Top 10 2021 và các nhóm tấn công thường gặp
ngoài Top 10. Tập luật gồm hai tầng: khoảng 45 \textit{chữ ký} (signature)
bắt mẫu tấn công hoàn chỉnh và 33 \textit{dấu hiệu} (indicator) là các tín
hiệu yếu cộng dồn theo mô hình anomaly scoring. Bảng \ref{table:rule-groups}
liệt kê các nhóm và số lượng rule trong mỗi nhóm. Danh sách chi tiết của 78
rule cùng pattern chính được trình bày ở Phụ lục B.

\begin{table}[h!]
\centering
\caption{Phân bố 78 rule theo 13 nhóm tấn công}
\label{table:rule-groups}
\small
\begin{tabular}{|l|c|p{6.5cm}|}
\hline
\textbf{Nhóm tấn công} & \textbf{Số rule} & \textbf{Đại diện pattern} \\
\hline
SQL Injection (SQLi) & 13 & UNION SELECT, OR 1=1, time-based sleep/benchmark, error-based extractvalue \\
\hline
Cross-Site Scripting (XSS) & 12 & \texttt{<script>}, event handler \texttt{onerror=}, \texttt{javascript:}/\texttt{data:} URI, iframe/svg \\
\hline
Remote Code Execution / CMDi & 16 & Shell metachar \texttt{; | \&\& \$()}, wget/curl/nc, PHP \texttt{system()}/\texttt{exec()}, reverse shell, JNDI \\
\hline
Path Traversal / LFI & 10 & \texttt{../}, biến thể mã hóa, \texttt{/etc/passwd}, \texttt{php://} wrapper, null byte \\
\hline
Server-Side Request Forgery (SSRF) & 6 & Dải IP nội bộ, \texttt{169.254.169.254}, \texttt{file://}, \texttt{gopher://} \\
\hline
XML External Entity (XXE) & 2 & \texttt{<!ENTITY}, SYSTEM, parameter entity \\
\hline
NoSQL Injection & 2 & Toán tử Mongo \texttt{\$where}, \texttt{\$ne}, \texttt{\$gt}, \texttt{\$regex} \\
\hline
Information Leak & 5 & Dotfiles/\texttt{.git}, file backup \texttt{.bak}/\texttt{.old}, file cấu hình, CRLF header injection \\
\hline
Scanner Detection & 4 & Scanner User-Agent (Nikto, sqlmap), probe path \texttt{.env}/\texttt{wp-admin} \\
\hline
Account Takeover (ATO) & 2 & Brute-force \texttt{/login}, lạm dụng password-reset \\
\hline
Bot / Automation & 1 & Thiếu \texttt{Accept-Language} (heuristic) \\
\hline
Denial of Service (DoS) & 1 & Cụm header quá khổ \\
\hline
Custom / Heuristic & 4 & Null byte, over-encoding, hex obfuscation, template injection (SSTI) \\
\hline
\end{tabular}
\end{table}

Cách xây dựng tập rule đi qua ba bước. Thứ nhất, lựa chọn các nhóm tấn
công có tần suất xuất hiện cao nhất trong các báo cáo bảo mật ứng dụng web
giai đoạn 2021--2024 (OWASP Top 10 2021, Verizon DBIR 2023, các CVE nổi
bật như Log4Shell). Thứ hai, với mỗi nhóm, khảo sát các pattern điển hình
từ OWASP CRS, từ payload list công khai (PayloadAllTheThings, SecLists),
và từ các sample tấn công đã được công bố. Thứ ba, viết và kiểm thử pattern
trên tập 42 OWASP attack payload đã biết, hiệu chỉnh transform chain và
proximity threshold cho đến khi đạt block rate \(\ge\) 95\%.

% ====================================================================
\section{DistilBERT -- Mô hình phân loại tấn công}
\label{section:3.3}

\subsection{Khái quát kiến trúc Transformer và BERT}
\label{subsection:3.3.1}

DistilBERT là một biến thể nhẹ của BERT, vốn là kiến trúc tiêu chuẩn cho
nhiều bài toán xử lý ngôn ngữ tự nhiên kể từ năm 2018. Để hiểu DistilBERT,
cần đi qua hai nền tảng tiền đề là kiến trúc Transformer và BERT.

\textbf{Transformer} (Vaswani và cộng sự, 2017) là kiến trúc mạng nơ-ron
được đề xuất ban đầu cho bài toán dịch máy, dựa hoàn toàn vào cơ chế
\textit{self-attention} thay vì các cell hồi quy (RNN/LSTM) hay convolution.
Cơ chế self-attention cho phép mỗi vị trí trong chuỗi đầu vào ``chú ý'' đến
mọi vị trí khác cùng một lúc, với trọng số được tính theo công thức
\(\text{softmax}\!\left(\frac{QK^\top}{\sqrt{d_k}}\right)V\) trong đó \(Q\),
\(K\), \(V\) là các ma trận query, key, value được tính từ chính chuỗi đầu
vào. Lợi thế của Transformer so với RNN nằm ở khả năng song song hóa hoàn
toàn quá trình tính toán cho mọi vị trí trong chuỗi, tận dụng tốt phần cứng
GPU/TPU hiện đại.

\textbf{BERT} (Devlin và cộng sự, 2018) là mô hình tiền huấn luyện
(pre-trained) dựa trên kiến trúc Transformer encoder hai chiều, gồm hai
nhiệm vụ tiền huấn luyện: \textit{Masked Language Modeling} (MLM -- dự đoán
các token bị che ngẫu nhiên trong văn bản) và \textit{Next Sentence
Prediction} (NSP -- dự đoán hai câu có liên tiếp hay không). Sau khi tiền
huấn luyện trên tập văn bản khổng lồ (Wikipedia + BookCorpus, khoảng 3.3 tỷ
từ), BERT có thể được fine-tune cho các nhiệm vụ downstream như phân loại
câu, gán nhãn token, hỏi đáp, với chi phí và lượng dữ liệu nhỏ hơn nhiều
so với huấn luyện từ đầu.

\textbf{DistilBERT} (Sanh và cộng sự, 2019) là phiên bản chắt lọc (distilled)
của BERT-base thông qua kỹ thuật \textit{knowledge distillation}: huấn
luyện một mô hình ``học sinh'' (student, kiến trúc đơn giản hơn) học cách
khớp với phân phối xác suất đầu ra của mô hình ``giáo viên'' (teacher, ở
đây là BERT-base). Kết quả là DistilBERT có 6 lớp encoder thay vì 12 (giảm
40\% số tham số xuống còn 66 triệu), nhanh hơn 60\% ở thời gian inference,
nhưng vẫn giữ được khoảng 97\% hiệu suất của BERT-base trên benchmark GLUE.
Đây chính là đặc tính khiến DistilBERT phù hợp với bài toán phân loại
real-time của WAF, nơi mỗi mili giây độ trễ đều ảnh hưởng đến trải nghiệm
người dùng cuối.

\subsection{Fine-tuning cho bài toán phân loại HTTP request}
\label{subsection:3.3.2}

Bài toán mà DistilBERT đảm nhận trong đề tài là \textit{sequence
classification} với 10 lớp: \texttt{normal} (không phải tấn công) và 9
lớp tấn công tương ứng với các nhóm chính (\texttt{sqli}, \texttt{xss},
\texttt{cmdi}, \texttt{path\_traversal}, \texttt{ssrf}, \texttt{xxe},
\texttt{log4shell}, \texttt{ssti}, \texttt{nosqli}). Trong giai đoạn
suy luận (inference), output cuối được giảm về phân loại nhị phân
(attack/normal) bằng cách gộp 9 nhãn tấn công thành một lớp -- phù hợp với
nhu cầu sử dụng của engine.

\textbf{Tiền xử lý đầu vào -- canonical text}: Đầu vào của mô hình là một
biểu diễn văn bản chuẩn hóa của HTTP request, gọi là \textit{canonical
text}. Đặc tả của canonical text được thống nhất giữa Go (hàm
\texttt{BuildCanonicalText} trong WAF core) và Python (hàm
\texttt{canonical\_compose} trong ML service) và phải tạo ra cùng một
chuỗi byte cho cùng một request -- đây là ràng buộc cứng được kiểm tra
bằng 7 fixture test case xuyên hai ngôn ngữ. Định dạng tổng quát:
\texttt{METHOD PATH\textbackslash nHost: ...\textbackslash nUser-Agent:
...\textbackslash n\textbackslash nBODY}, với các quy tắc cụ thể: URL-decode
một lớp đối với path và query, giới hạn mỗi header value ở 256 ký tự, sắp
xếp header theo thứ tự alphabetic để đảm bảo tính xác định.

\textbf{Pipeline fine-tuning}: Đề tài fine-tune từ checkpoint
\texttt{distilbert-base-uncased} với các tham số: optimizer AdamW, learning
rate \(2 \times 10^{-5}\), batch size 32, max sequence length 512 token,
số epoch tối đa 6 với early stopping theo validation F1. Loss function là
cross-entropy với class weights tỷ lệ nghịch với tần suất lớp để bù trừ
mất cân bằng dữ liệu. Dữ liệu huấn luyện được kiểm soát theo phiên bản:
v5 (baseline), v6 (fail report -- không đạt hard gate), v6.1 (kế hoạch
augmentation 1800 mẫu có mục tiêu).

\textbf{Hard gate kiểm soát chất lượng}: Đề tài áp đặt một bộ tiêu chí
cứng mà mô hình phải vượt qua trước khi được ship lên production: overall
accuracy \(\ge\) 0.92, cmdi F1 \(\ge\) 0.85, log4shell F1 \(\ge\) 0.90. Mô
hình không đạt gate sẽ không được tích hợp -- ví dụ v6 đã fail với accuracy
0.8866, dẫn đến quyết định giữ v5 trong production và lập kế hoạch v6.1.
Chi tiết của cơ chế hard gate được trình bày ở Chương \ref{section:5.3}.

\subsection{Lý do lựa chọn}
\label{subsection:3.3.3}

So với các kiến trúc khác có thể dùng cho cùng bài toán, DistilBERT cân
bằng tốt nhất giữa độ chính xác và tốc độ.

So với \textbf{LSTM hoặc CNN-text}, DistilBERT hưởng lợi từ tiền huấn luyện
trên kho ngữ liệu khổng lồ -- một lợi thế quan trọng khi dữ liệu HTTP request
có nhãn không đủ phong phú để huấn luyện một mô hình hoàn toàn từ đầu. Các
mẫu hình của BERT về cấu trúc ngữ pháp và mối quan hệ ngữ nghĩa đã được
học sẵn, fine-tuning chỉ phải điều chỉnh các phân biệt tinh tế giữa các
nhãn tấn công.

So với \textbf{BERT-base hoặc RoBERTa}, DistilBERT hi sinh khoảng 3\% accuracy
để đổi lấy gấp đôi tốc độ inference -- một đánh đổi mà đề tài chấp nhận
được, vì ML chỉ được gọi trong vùng xám (khoảng 5--10\% lưu lượng) và
timeout cứng 800 mili giây bao trùm tất cả các bước HTTP, tokenization, và
forward pass. Trên phần cứng CPU phổ thông (Intel i5 thế hệ 10), latency
inference của DistilBERT trung bình 60--80 mili giây cho một sequence ngắn,
phù hợp với budget thời gian.

So với các mô hình hệ thống học truyền thống như \textbf{Random Forest hay
SVM trên TF-IDF features}, DistilBERT vượt trội rõ rệt khi xử lý các biến
thể đối nghịch (adversarial variants) nhờ khả năng nắm bắt ngữ cảnh, không
phụ thuộc vào việc liệt kê các n-gram đặc trưng -- vốn dễ bị đánh lừa qua
character-level transformations.

% ====================================================================
\section{FastAPI và Python -- ML Inference Service}
\label{section:3.4}

\subsection{Khái quát}
\label{subsection:3.4.1}

FastAPI là framework xây dựng API HTTP cho Python được phát triển từ năm
2018, dựa trên ba thành phần nền tảng: \textbf{Starlette} (ASGI toolkit),
\textbf{Pydantic} (kiểm tra và validate dữ liệu qua type hint), và
\textbf{Uvicorn} (ASGI server hiệu năng cao dựa trên uvloop và httptools).
Đặc điểm nổi bật của FastAPI bao gồm: tự động sinh tài liệu OpenAPI (Swagger
UI) từ type hint, tích hợp sẵn async/await cho I/O concurrency, và hiệu
năng cao hơn các framework cổ điển Flask/Django đáng kể (ở mức gần ngang
với Node.js Express trong các benchmark độc lập).

\subsection{Vai trò trong hệ thống}
\label{subsection:3.4.2}

ML Inference Service đóng vai trò microservice chuyên trách suy luận, được
viết bằng Python để tận dụng hệ sinh thái Hugging Face (transformers,
tokenizers) và PyTorch -- hai gói không có tương đương trưởng thành trong
hệ sinh thái Go. Service cung cấp hai endpoint chính: \texttt{POST /predict}
nhận một canonical text và trả về \texttt{(label, confidence)}; và
\texttt{POST /predict\_batch} nhận một mảng canonical text và trả về mảng
kết quả tương ứng -- được sử dụng cho mục đích đánh giá hàng loạt.

Khi khởi động, service nạp checkpoint model từ thư mục local
(\texttt{model\_v5/}, \texttt{model\_v6/}), chuyển sang chế độ \texttt{eval}
của PyTorch, và (khi cấu hình) bật chế độ inference no-grad để giảm tiêu
thụ bộ nhớ. Trên CPU, một worker đơn xử lý tuần tự các request đến; khi cần
mở rộng, có thể tăng số worker Uvicorn hoặc chạy nhiều instance phía sau
một load balancer.

Việc tách ML thành service riêng có ba lợi ích cụ thể: \textbf{(i)} cách
ly môi trường runtime, không trộn lẫn binary Go với phụ thuộc Python; \textbf{(ii)}
có thể nâng cấp hoặc thay thế mô hình (đổi từ DistilBERT sang một mô hình
khác) mà không cần biên dịch lại WAF core; và \textbf{(iii)} cho phép chia
sẻ service ML giữa nhiều instance WAF trong tương lai, tận dụng cache và
warm-up của mô hình.

% ====================================================================
\section{PostgreSQL -- Lưu trữ logs, users, cấu hình}
\label{section:3.5}

\subsection{Khái quát}
\label{subsection:3.5.1}

PostgreSQL là hệ quản trị cơ sở dữ liệu quan hệ mã nguồn mở, phát triển
liên tục từ năm 1996, được biết đến với độ tuân thủ chuẩn SQL cao, hỗ trợ
phong phú các kiểu dữ liệu (JSON, JSONB, mảng, range, network), và khả
năng mở rộng qua extension (PostGIS, pg\_trgm, TimescaleDB). Phiên bản đề
tài sử dụng là PostgreSQL 15, hỗ trợ đầy đủ \texttt{CHECK constraint},
\texttt{GENERATED column}, và partitioning theo timestamp -- ba tính năng
hữu ích cho schema của đề tài.

\subsection{Vai trò trong hệ thống}
\label{subsection:3.5.2}

PostgreSQL đảm nhận ba trách nhiệm chính: \textbf{(i)} lưu trữ tài khoản
người dùng nội bộ (bảng \texttt{users}), \textbf{(ii)} ghi nhật ký kiểm
toán (bảng \texttt{audit\_logs}), và \textbf{(iii)} lưu trữ động các danh
sách IP (bảng \texttt{ip\_lists}). Cấu hình WAF cốt lõi (ngưỡng, URL backend)
được lưu trong tệp \texttt{config.yaml} thay vì database để giảm phụ thuộc
khởi động.

Việc lựa chọn PostgreSQL thay vì các SQL khác (MySQL, SQLite) hay NoSQL
(MongoDB) dựa trên hai cân nhắc. Thứ nhất, dữ liệu audit log có tính chất
quan hệ rõ ràng và yêu cầu truy vấn lọc/sắp xếp/aggregate phức tạp (theo
thời gian, IP, ruleID, kiểu tấn công) -- phù hợp với mô hình quan hệ và
chỉ mục B-tree truyền thống. Thứ hai, hỗ trợ \texttt{CHECK constraint} và
\texttt{UNIQUE} ràng buộc ở mức schema là cốt lõi của chính sách bảo mật
nội bộ (chống nhập role bất hợp lệ, chống trùng username/email -- mục
\ref{subsection:2.2.8}).

\subsection{Schema chính}
\label{subsection:3.5.3}

\textbf{Bảng \texttt{users}} lưu thông tin tài khoản nội bộ. Các cột:
\texttt{id} (\texttt{SERIAL PRIMARY KEY}), \texttt{username}
(\texttt{VARCHAR UNIQUE NOT NULL}), \texttt{email} (\texttt{VARCHAR UNIQUE}),
\texttt{password\_hash} (\texttt{VARCHAR NOT NULL} -- bcrypt hash, mục
\ref{section:3.6}), \texttt{role} (\texttt{VARCHAR CHECK (role IN ('admin',
'editor', 'viewer'))}), \texttt{created\_at}, \texttt{updated\_at}.

\textbf{Bảng \texttt{audit\_logs}} lưu mỗi bản ghi kiểm toán. Các cột chính:
\texttt{id}, \texttt{timestamp} (có index), \texttt{ip}, \texttt{method},
\texttt{path}, \texttt{score}, \texttt{decision} (ENUM \texttt{ALLOW}/\texttt{LOG}/\texttt{CHALLENGE}/\texttt{BLOCK}),
\texttt{rule\_ids} (\texttt{INTEGER[]}), \texttt{ml\_label}, \texttt{ml\_confidence},
\texttt{actor\_id} (\texttt{NULLABLE}, tham chiếu \texttt{users.id} cho các
sự kiện hệ thống), \texttt{event\_type} (request hoặc system).

\textbf{Bảng \texttt{ip\_lists}} lưu các IP/CIDR thuộc blacklist hoặc
whitelist. Các cột: \texttt{id}, \texttt{ip\_or\_cidr}, \texttt{list\_type}
(\texttt{CHECK IN ('blacklist','whitelist')}), \texttt{expires\_at}
(\texttt{NULLABLE} -- hỗ trợ TTL tự động hết hạn), \texttt{reason}.

% ====================================================================
\section{JWT và Bcrypt -- Xác thực và bảo mật}
\label{section:3.6}

\subsection{JWT: xác thực không trạng thái với dual-mode delivery}
\label{subsection:3.6.1}

JSON Web Token (JWT) là chuẩn được mô tả trong RFC 7519, dùng để truyền
khẳng định (claim) giữa các bên dưới dạng một chuỗi ba phần
\texttt{header.payload.signature} được mã hóa base64url. Đề tài sử dụng
thuật toán ký \texttt{HS256} (HMAC-SHA256 với khóa bí mật) -- phù hợp cho
mô hình một service đơn lẻ xác thực và tiêu thụ token cho chính nó. Trong
các kiến trúc phức tạp hơn với nhiều service tiêu thụ, \texttt{RS256} với
cặp khóa bất đối xứng là lựa chọn tốt hơn, nhưng nằm ngoài phạm vi của đề
tài.

Đặc tính cốt lõi của JWT là tính \textit{stateless} -- server không cần
lưu trạng thái session ở phía mình; chỉ cần verify chữ ký và đọc claim.
Đặc tính này phù hợp với mô hình triển khai container hóa của đề tài: các
instance WAF có thể được scale ngang mà không cần share session store.

\textbf{Dual-mode delivery}: Đề tài hiện thực hóa hai cơ chế song song để
client gửi JWT về server. Cơ chế thứ nhất là \textit{HttpOnly cookie} --
JWT được set vào cookie với cờ \texttt{HttpOnly} (JavaScript không thể đọc,
chống XSS exfiltration), \texttt{Secure} (chỉ gửi qua HTTPS), và
\texttt{SameSite=Lax} (giảm thiểu CSRF). Đây là cơ chế mặc định cho dashboard
web. Cơ chế thứ hai là \textit{Authorization Bearer header} -- client gửi
header \texttt{Authorization: Bearer <token>}; phù hợp cho API client và
các kịch bản scripting. Middleware xác thực kiểm tra cả hai nguồn, cho phép
một backend thống nhất phục vụ cả browser-based dashboard lẫn API
integration.

\subsection{Bcrypt: hashing mật khẩu với cost factor}
\label{subsection:3.6.2}

Bcrypt là hàm hash mật khẩu được Niels Provos và David Mazières giới thiệu
năm 1999, dựa trên thuật toán mã hóa Blowfish nhưng được điều chỉnh để
trở thành hàm hash một chiều, cố ý chậm. Đặc tính then chốt của bcrypt là
tham số \textit{cost factor} -- số mũ trên một biến đếm vòng lặp nội bộ
(work factor = \(2^{\text{cost}}\)) -- cho phép tăng chi phí tính toán
theo thời gian khi phần cứng nhanh hơn, mà không thay đổi giao diện hàm.

Đề tài sử dụng cost factor 10 (1024 vòng) -- giá trị mặc định của gói
\texttt{golang.org/x/crypto/bcrypt}, tương ứng với khoảng 60--80 mili giây
để hash một mật khẩu trên phần cứng phổ thông. Mức cost này được khuyến nghị
bởi OWASP Password Storage Cheat Sheet (2023) như mức cân bằng giữa an toàn
chống brute force và trải nghiệm người dùng (đăng nhập không bị cảm nhận
được độ trễ).

So với các hàm hash mật khẩu hiện đại hơn (Argon2, scrypt), bcrypt vẫn được
coi là chấp nhận được cho hầu hết các use case và có lợi thế là thư viện
sẵn sàng, được kiểm chứng, và không yêu cầu cấu hình bộ nhớ phức tạp. Đề
tài chấp nhận đánh đổi này do tập user nội bộ có quy mô nhỏ, không phải
target chính của các cuộc tấn công brute force quy mô lớn.

\subsection{RBAC: ba vai trò admin / editor / viewer}
\label{subsection:3.6.3}

Role-Based Access Control (RBAC) là mô hình phân quyền cổ điển trong đó
quyền truy cập được gán cho \textit{vai trò} thay vì cho cá nhân; người dùng
được gán một (hoặc nhiều) vai trò, và quyền của họ là tổng hợp các quyền
của các vai trò đó. Đề tài hiện thực hóa phiên bản đơn giản nhất của RBAC
với mỗi người dùng có đúng một vai trò.

\textbf{Ma trận phân quyền}: Bảng \ref{table:rbac-matrix} liệt kê các thao
tác nhạy cảm và vai trò có quyền thực hiện.

\begin{table}[h!]
\centering
\caption{Ma trận phân quyền RBAC ba vai trò}
\label{table:rbac-matrix}
\small
\begin{tabular}{|p{6cm}|c|c|c|}
\hline
\textbf{Thao tác} & \textbf{admin} & \textbf{editor} & \textbf{viewer} \\
\hline
Xem dashboard, audit log & \checkmark & \checkmark & \checkmark \\
\hline
Cập nhật thông tin cá nhân (\texttt{/me}) & \checkmark & \checkmark & \checkmark \\
\hline
Đổi mật khẩu của mình (\texttt{/me/password}) & \checkmark & \checkmark & \checkmark \\
\hline
Sửa rule, IP list, system settings & \checkmark & \checkmark & -- \\
\hline
Cấu hình kênh cảnh báo & \checkmark & \checkmark & -- \\
\hline
CRUD tài khoản người dùng khác & \checkmark & -- & -- \\
\hline
Reset mật khẩu của người khác & \checkmark & -- & -- \\
\hline
\end{tabular}
\end{table}

\textbf{Hiện thực hóa}: Vai trò được lưu ở cột \texttt{role} của bảng
\texttt{users} với ràng buộc \texttt{CHECK} ở mức schema (mục
\ref{subsection:3.5.3}). Tại runtime, hai middleware được sử dụng:
\texttt{requireAuthN} chỉ yêu cầu JWT hợp lệ (cho phép mọi role có token),
và \texttt{requireAdmin} thêm điều kiện \texttt{claim.role = "admin"}.
Các endpoint admin-only được gate qua \texttt{requireAdmin}; các endpoint
chỉ yêu cầu đăng nhập (như \texttt{/me}) qua \texttt{requireAuthN}.

\textbf{Bảo vệ chống lockout}: Đặc thù của bài toán quản trị tài khoản là
khả năng tự đẩy hệ thống vào trạng thái không thể phục hồi -- ví dụ admin
duy nhất tự xóa hoặc tự demote chính mình, dẫn đến hệ thống không còn admin
nào. Để chống lại, đề tài định nghĩa hai bất biến (invariant) được kiểm tra
ở tầng handler trước khi commit thay đổi: \textit{last-admin protection}
(không cho phép xóa hoặc demote admin cuối cùng, dựa trên truy vấn
\texttt{SELECT COUNT(*) FROM users WHERE role = 'admin'}) và \textit{self-action
protection} (admin không được xóa/demote chính tài khoản đang đăng nhập).
Hai bất biến này được trình bày chi tiết ở mục \ref{subsection:2.2.10}.

% ====================================================================
\section{Token Bucket -- Thuật toán Rate Limiting}
\label{section:3.7}

\subsection{Khái quát Token Bucket}
\label{subsection:3.7.1}

Token Bucket là một thuật toán điều khiển tải đường truyền cổ điển được
mô tả lần đầu trong bối cảnh quản lý chất lượng dịch vụ mạng (QoS) cuối
thập niên 1980. Mô hình hoạt động dựa trên hình ảnh trực quan của một
\textit{bucket} (xô) chứa các \textit{token} (thẻ). Bucket có hai tham số:
\textit{rate} (số token được rót vào bucket mỗi giây) và \textit{burst}
(kích thước tối đa của bucket). Mỗi request đến cần lấy đi một token mới
được phục vụ; nếu bucket trống, request bị từ chối hoặc được đưa vào hàng
chờ.

Đặc tính then chốt của Token Bucket là cho phép \textit{burst} ngắn hạn:
nếu một thời gian dài không có request, bucket được tích lũy đầy token và
có thể phục vụ một spike traffic mà không phạt; sau đó traffic phải trở về
dưới rate trung bình để bucket có thời gian phục hồi. Hành vi này phù hợp
hơn với traffic web thực tế (có tính burst tự nhiên) so với mô hình
\textit{leaky bucket} cổ điển vốn áp đặt rate cứng tại mọi thời điểm.

\subsection{So sánh với các thuật toán khác}
\label{subsection:3.7.2}

So với \textbf{Sliding Window Counter}, Token Bucket nhẹ hơn về bộ nhớ
(chỉ cần lưu hai giá trị: số token hiện tại và timestamp lần cập nhật),
trong khi Sliding Window cần lưu chuỗi timestamp của các request gần đây.
Tuy nhiên, Sliding Window cho điểm chính xác hơn ở biên cửa sổ -- một điểm
yếu của Token Bucket ít quan trọng trong bài toán của đề tài.

So với \textbf{Fixed Window Counter}, Token Bucket không gặp vấn đề
``boundary spike'' -- khi attacker gửi burst kép ở hai bên ranh giới cửa sổ
để vượt qua giới hạn danh nghĩa. Mỗi bucket trong Token Bucket cập nhật
liên tục theo thời gian, không có ranh giới rạch ròi.

\subsection{Tích hợp vào middleware stack}
\label{subsection:3.7.3}

Trong hệ thống, mỗi địa chỉ IP có một bucket riêng. Toàn bộ tập bucket được
lưu trong một \texttt{sync.Map} cho phép truy cập concurrent không cần lock
tường minh. Một goroutine quét định kỳ (mỗi 5 phút) loại bỏ các bucket có
trạng thái không hoạt động lâu hơn 10 phút, ngăn chặn rò rỉ bộ nhớ khi
IP space biến động lớn.

Rate Limiter được đặt trong middleware stack sau \textit{IP-check middleware}
(blacklist/whitelist) và trước \textit{rule engine middleware}, theo nguyên
tắc \textit{cheap check first}: từ chối sớm các request không cần xử lý sâu.
Khi một IP vượt quá ngưỡng cho phép liên tục, IP đó được tự động bổ sung
vào blacklist động với TTL ngắn (mặc định 1 giờ) -- đây là cơ chế phòng vệ
hành vi (behavior-based defense) bổ sung cho rule pattern matching.

% ====================================================================
\section{Docker và Docker Compose -- Đóng gói và triển khai}
\label{section:3.8}

\subsection{Containerization các thành phần}
\label{subsection:3.8.1}

Docker là công nghệ container hóa được phát triển từ năm 2013, đóng gói
ứng dụng cùng toàn bộ phụ thuộc runtime vào một image duy nhất, đảm bảo
tính tái lập của môi trường thực thi. Đề tài đóng gói ba thành phần riêng
biệt: \textbf{(i)} WAF core (binary Go biên dịch tĩnh, image dựa trên
\texttt{scratch} hoặc \texttt{distroless} với kích thước dưới 30 MB);
\textbf{(ii)} ML Inference Service (image dựa trên \texttt{python:3.13-slim}
kèm phụ thuộc PyTorch và transformers, kích thước khoảng 1.5 GB do model
weight); và \textbf{(iii)} PostgreSQL (image chính thức \texttt{postgres:15}).

Việc đóng gói theo Docker mang lại ba lợi ích cụ thể. Thứ nhất, tách biệt
runtime: WAF binary không cần Python hay PyTorch ở máy đích, và ngược lại.
Thứ hai, cùng image có thể chạy trên mọi nền tảng host hỗ trợ Docker, từ
máy trạm phát triển đến server sản xuất, không phụ thuộc hệ điều hành cụ
thể. Thứ ba, image có thể được publish lên registry (Docker Hub, GHCR) và
phân phối với version tag rõ ràng, hỗ trợ rollback nhanh khi cần.

\subsection{Docker Compose one-command deployment}
\label{subsection:3.8.2}

Docker Compose là công cụ điều phối multi-container ở quy mô một máy
chủ đơn, dựa trên tệp YAML mô tả tập service, mạng và volume. Đề tài cung
cấp tệp \texttt{docker-compose.db.yml} mô tả ba service nói trên, với các
cấu hình quan trọng: ánh xạ cổng (\texttt{8080}, \texttt{8443} cho WAF;
\texttt{8000} cho ML; \texttt{5432} cho PostgreSQL), volume mount cho dữ
liệu PostgreSQL (\texttt{postgres\_data} -- persistent), và biến môi trường
nạp từ tệp \texttt{.env}.

Toàn bộ triển khai được khởi động bằng một lệnh: \texttt{docker compose up
-d}. Đối với người dùng mới, đề tài cung cấp thêm script
\texttt{scripts/setup\_db.sh} thực hiện migration ban đầu (tạo schema, seed
admin mặc định) và \texttt{scripts/generate\_certs.sh} sinh self-signed TLS
certificate cho môi trường thử nghiệm.

\subsection{Lý do lựa chọn}
\label{subsection:3.8.3}

So với các giải pháp đóng gói khác như \textbf{tar archive thuần} (yêu cầu
cài thủ công phụ thuộc) hay \textbf{Kubernetes manifest} (phức tạp quá mức
cho quy mô đề tài), Docker Compose ở mức cân bằng phù hợp: đủ mạnh để mô tả
multi-service deployment với network và volume rõ ràng, đủ đơn giản để một
sinh viên hay kỹ sư DevOps mới có thể vận hành mà không cần hạ tầng cluster.

Đối với các triển khai sản xuất quy mô lớn hơn, Docker Compose có thể được
chuyển đổi sang Kubernetes manifest (qua công cụ \texttt{kompose}) hoặc
Helm chart, nhưng phép thay đổi này không thuộc phạm vi đề tài.

\bigskip
\noindent\textbf{Kết chương.}
Chương này đã trình bày tám lớp công nghệ cốt lõi tạo nên hệ thống: Go cho
WAF core với mô hình goroutine và GC độ trễ thấp; Rule Engine với mô hình
anomaly scoring và transform chain chống evasion; DistilBERT cho lớp ML
xác minh vùng xám; FastAPI/Python cho ML Inference Service tách biệt
runtime; PostgreSQL cho lưu trữ quan hệ với ràng buộc tính toàn vẹn ở mức
schema; JWT và bcrypt cho xác thực không trạng thái và lưu mật khẩu an
toàn, kết hợp với RBAC ba vai trò admin/editor/viewer cùng các bất biến
bảo vệ chống lockout; Token Bucket cho điều khiển tải per-IP với khả năng
chịu burst; và Docker Compose cho đóng gói triển khai một lệnh. Mỗi lựa
chọn được lý giải thông qua so sánh với các ứng cử viên thay thế và đối
chiếu với các yêu cầu phi chức năng đã xác định ở mục \ref{section:2.4}.
Đây là cơ sở công nghệ trực tiếp cho phần thiết kế và hiện thực hóa hệ
thống được trình bày trong Chương \ref{section:4.1}.

\end{document}
